% Reciprocal proof completion for ED26-REQ-N-c-004-03.

\section{Immerman's ordered fixed-point capture}
\label{sec:annex-p09-n-c-004-03}

\begin{annexitem}{annex:p9.n-c-004.3.evaluation}{thm:p09-immerman1986-evaluation}{Positivity, the iteration bound, and the fixed input axis}
\begin{proof}
Put a first-order matrix in negation-normal form.  If \(R\) is positive,
enlarging its interpretation preserves truth at its positive atoms, and
conjunction, disjunction and both quantifiers preserve this implication.
The other atoms are unchanged.  To cover nested fixed points and
intervening negations, prove both polarities together by induction on
syntax: a positive free relation parameter is covariant, and a negative
one is contravariant.  Negation interchanges these two implications;
the Boolean operations and quantifiers preserve them.  At an inner
\(\mu S\) constructor, its operator is monotone in its own bound
\(S\).  If the free parameter \(R\) is positive in the defining
matrix and \(R\subseteq R'\), the syntax induction gives
\(F_R(A)\subseteq F_{R'}(A)\) for every candidate \(A\).
Starting both iterations at the empty relation, induction gives
\(F_R^i(\varnothing)\subseteq F_{R'}^i(\varnothing)\): use
pointwise operator inclusion and then monotonicity in \(S\).
Their stable unions have the same inclusion.  For a negative parameter
the pointwise operator inclusion and the stage inclusion are reversed.
Consequently the binder preserves the matrix's parameter polarity, and
any surrounding negation flips it.  This also covers a parameter
negative inside an inner induction but positive after a surrounding
negation.  Thus the admitted formula is monotone in its bound relation.

From \(I_0=\varnothing\subseteq I_1\) and monotonicity, induction
gives \(I_i\subseteq I_{i+1}\).  This ascending chain does not require
the operator to be inflationary on \emph{every} arbitrary relation.
If two successive stages are equal, applying the same operator makes
all later stages equal.  Otherwise at least one new tuple enters.
There are only \(n^r\) tuples, so the chain stabilizes after at
most \(n^r\) strict increases.  Its stable value \(I\) is a
fixed point.  If \(J\) is any fixed point, induction gives
\(I_i\subseteq J\), starting from the empty set and applying
monotonicity; hence \(I\subseteq J\).  It is the least fixed point.

For polynomial evaluation, induct on the fixed query's syntax.  Atomic
lookups and truth tables cost polynomial time on explicit relation tables.
Quantifiers enumerate at most \(n\) values per variable.  At a
fixed-point constructor, at most \(n^r\) stages each evaluate the
fixed defining formula on at most \(n^r\) candidate tuples; the
syntax induction bounds each such evaluation polynomially, including its
nested fixed points and fixed parameter tuple.  Products of these finitely
many fixed polynomials remain polynomial.  Enumerating the entire output
relation costs a further factor \(n^d\), with \(d\) fixed.
The exponent depends on the query's variable counts, arities and nesting;
letting the formula vary does not preserve this fixed exponent.

For the family example, induction on stages gives paths of length at
most \(i-1\) at stage \(i\geq1\): the diagonal supplies length
zero, and prepending one parent edge supplies the next length bound.
Conversely every generated pair has such a path by the same induction.
The stable value is therefore exactly reflexive transitive closure,
not strict ancestry.  This completes the source's iteration argument
and fixes the meaning of its displayed example \cite[secs.~1--2]{Immerman1986}.
\end{proof}
\end{annexitem}

\begin{annexitem}{annex:p9.n-c-004.3.initial}{thm:p09-immerman1986-initial}{The printed binary row and the fixed-type characteristic-table extension}
\begin{proof}
In the source's binary case the first clause of \(M_0\) applies only
at the first position and picks code \(4\) if \(E(0,0)\) holds
and code \(3\) otherwise.  The second applies exactly to the other
positions with \(p_1=0\), and uniquely picks code \(1\) or
\(0\) according to \(E(p_2,p_3)\).  The third applies at
all remaining positions and uniquely supplies blank.  These position
cases are disjoint and exhaustive.  There are \(n^2\) positions
with \(p_1=0\), so the stated bit-block and blank-block lengths
follow.  Each numbered code can be defined by a first-order formula
saying that it is the corresponding fixed-rank element of the order:
the least element has rank zero, and a fixed finite chain of immediate
successors defines each subsequent code.  These are definable elements,
not new input constants.  Each displayed equality with a code is eliminated
by the corresponding first-order code predicate.  If the code rank
does not exist, that predicate is false, rather than denoting a partial
term.  Their existence for the working alphabet needs the stated cutoff.

Here is a fully specified generalization of the source's table-decoding
step.  It is an editorial padded-block completion of Lemma~3.1, not a
claim that the paper printed this layout.  Make a fixed list of tables:
all \(R_i\), the available order if it is part of the query's input,
and the unary singleton tables \(x=c_i\).  Let \(r\geq1\)
be at least every table arity and choose \(k\geq r+2\).
Reserve finitely many distinct fixed-rank position elements as block
tags.  Use a \(k\)-tuple address of the form
\[
 (\underbrace{0,\ldots,0}_{k-r-1},b,u_1,\ldots,u_r).
\]
For a table \(T\) of arity \(a\leq r\) at tag \(b\), its
active positions have \(u_1=\cdots=u_{r-a}=0\), and its bit is
the truth value of \(T(u_{r-a+1},\ldots,u_r)\).
This is a finite conjunction of coordinate equalities and one input
atom, or its negation for bit zero.  A singleton table uses the atom
\(u_r=c_i\).  Every other position in that block is blank.
Different tags make the active blocks disjoint.  Positions beyond the
finitely many tags, or outside the displayed zero prefix, are also blank.
Thus all letter cases are first-order and exhaustive.

To make the domain size decodable even for low-arity or all-empty input
tables, reserve tag zero for a block whose first \(n\) positions
are bit one and whose remaining positions are blank: the active condition
is \(u_1=\cdots=u_{r-1}=0\), with \(u_r\) unrestricted.
Reserve tag one entirely blank.  Hence the first run of nonblank input
bits has length exactly \(n\), including when \(r=1\).  All
actual tables use tags starting at two.  This is an explicit unary-size
header for the completed layout, not an assumption that a machine can
infer the universe from a relation's positive tuples.  From that header
the machine computes the fixed block addresses and recovers every dense
table and constant value in polynomial time.  Conversely, constructing
this padded encoding from a dense explicit database takes polynomial
time.  The active blocks fit within \(O(n^r)\) tape positions since
their number is fixed.  Padding to \(n^k\) leaves polynomial room
for the computation.  The coding is not exponentially succinct.

At the first address replace its ordinary initial bit by the matching
start-state/head letter, exactly as the source replaces its first \(E\)
bit by code \(3\) or \(4\).  All other bits remain ordinary
letters.  The finite alphabet's alternative coding, including a specified
\(c_0\) code for acceptance, merely replaces the finite letter-equality
tests by the corresponding definable code tests.  There is still exactly
one initial letter at each position and exactly one head marker.
For each finite output arity, additional tested-tuple coordinates are
represented by singleton tables \(u_r=x_i\); these are first-order
parameter atoms, not extra quantified certificate relations.

All tags, letter tests, arities and coordinate constraints are fixed by
the machine and type.  A fixed larger cutoff provides the finitely many
distinct elements needed for them.  Databases below this cutoff have
only finitely many ordered isomorphism types and will be handled by
finite diagrams in the capture proof.  This supplies the general-type
claim that follows the source's printed binary formula without appealing
to a different descriptive-complexity theorem \cite[lem.~3.1]{Immerman1986}.
\end{proof}
\end{annexitem}

\begin{annexitem}{annex:p9.n-c-004.3.tableau}{thm:p09-immerman1986-tableau}{Local deterministic transitions and exactness of every induction row}
\begin{proof}
A deterministic one-tape transition \((q,b)\mapsto(q',b',d)\)
has \(d\in\{-1,0,1\}\).  At the current head cell it writes
\(b'\), carrying marker \(q'\) there if \(d=0\).
If the head moves right, the cell to its right gains marker \(q'\)
on that cell's previous tape symbol; if it moves left, the cell to its
left gains the marker.  Every other cell retains its previous ordinary
symbol.  Thus the new letter at a cell is determined by the old letters
at it and its two neighbors.  Inspect the finite set of triples to
produce the finite table \(\delta_M\).  On triples consistent with
a legal configuration its output is unique.  The normalized machine's
absorbing accepting and rejecting configurations use the stationary
rules.  Choices for impossible multiple-head triples are immaterial,
because the induction below never generates them.

At position zero the left neighbor is \(\#\) and the machine's
fixed left-end rules prohibit moving left of zero.  At the last position
the right neighbor is \(\#\); by the time/space normalization the
head never reaches that position.  The last ordinary cells therefore
remain blank, and the local table has the corresponding unchanged-blank
rules.  These clauses define the endpoints rather than silently applying
an undefined \(\bar p-1\) or \(\bar p+1\).  Every interior
cell has exactly one immediate predecessor and successor in the
lexicographic position order, as given by \(\operatorname{Succ}_k^p\).
Every nonfirst time has exactly one predecessor in the time order.

Let \(J_i\) be the \(i\)-th iterate of \(\varphi_M\) from
the empty relation.  Prove that \(J_i\), for \(i\leq n^k\),
contains exactly all tape-letter facts of times \(0,\ldots,i-1\).
For \(i=0\) it is empty.  At the first application, the base clause
supplies precisely the initial row by \(M_0\); the successor
clause cannot use a center-cell \(R\)-fact, since none exists.
Suppose the claim holds at stage \(i\).  The base clause regenerates
the initial row.  For time \(t>0\), the successor clause requires
the complete, uniquely coded neighborhood at time \(t-1\), with
the indicated exterior letter where needed.  If \(t-1<i\), all
these facts are available, the unique local table entry gives exactly
the correct new letter, and no other letter can be generated.
If \(t-1\geq i\), the center-cell fact is missing, so no
fact for time \(t\) is generated.  This proves that stage
\(i+1\) contains exactly rows through time \(i\), with one
letter per cell and no extras.  It also proves persistence of earlier
rows without adding an inflationary \(R\lor\varphi_M\) clause
that the source does not use.

After \(n^k\) applications all times are present.  The next application
regenerates precisely those rows: there is no time successor beyond the
last time.  The relation is a fixed point, and positivity plus iteration
from the empty set makes it the least one.  The final time lies after
halting by normalization; its stationary first-position letter is
\(q_f\) exactly for an accepting input.  This establishes the actual
cell/time/letter readout of the source's section~3 \cite[thm.~1]{Immerman1986}.
\end{proof}
\end{annexitem}

\begin{annexitem}{annex:p9.n-c-004.3.capture}{thm:p09-immerman1986-capture}{The source's constant anchor, finite exceptions, and return to queries}
\begin{proof}
The evaluation lemma proves \(\mathcal L_\mu(\leq)\subseteq
\mathrm{QPTIME}\), since lookup of a given order relation is
polynomial on the same explicit inputs.  A one-fixed-point formula is
an \(\mathcal L_\mu\) formula, so it remains to construct that
normal form for every polynomial-time query.

Fix a Boolean query and its deterministic polynomial-time machine on
explicit database encodings.  A standard one-tape simulation has only
polynomial slowdown: encode the finitely many old tapes as delimited
tracks on one tape, scan the at-most-polynomially-many visited cells
to locate their head markers and symbols, and scan again to perform
one old transition.  A polynomial number of transitions times a
polynomial scan length remains polynomial.  In an initial setup step,
the head creates a permanent origin flag alongside the first cell's
data; it knows that cell because it starts there.  The flag is a finite
alphabet component, its data are preserved under subsequent writes,
and the simulation can then recognize the first cell on return.  The
initial \(M_0\) still has the source's start-state letter at that
cell; this flag is created by a transition, not by silently changing
the printed initial row.  Returning to the first cell after acceptance
or rejection costs at most the visited tape length.
Writing its designated terminal letter and then remaining stationary
completes the normalization.  The decoder from the initial-row lemma
has polynomial cost as well.  Choose \(k\) larger than these fixed
degrees and enlarge the fixed cutoff so that their total time and space
are below \(n^k-2\).  No resources depend on a formula supplied
later as input.

Here are explicit orders implementing the source's footnote~3.  From
the given strict order \(<\), define the position order
\[
 x<_p y\ \Longleftrightarrow\
 (x=c_0\land y\ne c_0)\lor
 (x\ne c_0\land y\ne c_0\land x<y),
\]
and the time order
\[
 x<_t y\ \Longleftrightarrow\
 (x\ne c_0\land y=c_0)\lor
 (x\ne c_0\land y\ne c_0\land x<y).
\]
Both are total orders: they move the distinguished element to the first
or last place and leave the other elements in their original order.
Their tuple orders, successors, and extrema are first-order by substitution
in the formulas already displayed.  Thus the first position is
\(\bar c_0\) and the last time is \(\bar c_0\), while their
successor meanings are different.  Assign the accepting-letter code to
\(c_0\) and assign every other member of the fixed alphabet to a
distinct fixed-rank non-\(c_0\) element in the original order.
These code tests are also first-order for large enough domains.

This recoding does not discard the input order: its characteristic
table is available to the decoder for an order-sensitive query.
A machine can recover the original order and recode the database to
its chosen dense representation in polynomial time.  Equivalently it
can compute the fixed query directly from all decoded relations.
The position/time orders only determine tableau addresses; they do
not replace the distinguished input order in the query being decided.
The initial-row lemma and the tableau lemma now give
\[
 (\bar c_0,\bar c_0,c_0:\mu R(\bar p,\bar t,a))\varphi_M
 \quad\Longleftrightarrow\quad M\hbox{ accepts }B
\]
on all databases above the fixed cutoff.  This is exactly one
fixed-point expression with a first-order positive defining matrix.

For the omitted small sizes, enumerate their finitely many ordered
isomorphism types that satisfy the fixed query.  A complete diagram
for a type of size \(m\) existentially names \(m\) distinct
points, asserts that every point equals one of them, states the truth
or falsity of each relation on each tuple of them, and states which
named point interprets each constant.  The disjunction of the accepting
diagrams is a first-order sentence \(\operatorname{SmallYes}\).
Let \(\operatorname{Large}\) say there are more points than the
fixed cutoff.  Replace the defining matrix by
\[
\begin{split}
 \widehat\varphi_M(R)[\bar p,\bar t,a]\defeq{}&
  (\operatorname{SmallYes}\land\bar p=\bar c_0
     \land\bar t=\bar c_0\land a=c_0)\\
 &{}\lor(\operatorname{Large}\land\varphi_M(R)[\bar p,\bar t,a]).
\end{split}
\]
It remains positive in \(R\).  On a small accepting database it
generates the readout tuple in its first stage; on a small rejecting
database it generates nothing; on a large database it is the actual
tableau operator.  Thus the single readout expression works at every
positive size, including one-element structures.  The diagram sentence
does not pretend to supply missing distinct alphabet codes.

For a fixed \(d\)-ary query \(Q\), run the membership machine on
\((B,\bar x)\) and include the unary singleton tables \(u=x_i\)
in the decoder, as in the initial-row lemma.  The parameters \(\bar x\)
remain free in the operator; they are not part of the recursively bound
tuple and do not change during iteration.  Complete small diagrams also
record those parameters.  The same one-fixed-point expression then
defines membership for every \(\bar x\).  This proves the return
from the source's Boolean simplification to the asserted fixed-arity
query equality.

If no constant exists, use one free variable \(z\) as the anchor in
the position, time and alphabet coding.  The resulting readout has the
same truth value for every choice of \(z\): it decodes the same
database with a different address/code bijection.  Nonemptiness provides
an anchor, and one final \(\exists z\) closes the Boolean query.
This is the stated source modification, not a claim that a closed
positive-arity fixed-point atom has ground arguments without terms.
We have proved \(\mathrm{QPTIME}\subseteq\Sigma_\omega(\leq)\)
with the appropriate constant proviso and completed both inclusions
\cite[thm.~1 and its footnote~3]{Immerman1986}.
\end{proof}
\end{annexitem}

\begin{annexitem}{annex:p9.n-c-004.3.unordered}{thm:p09-immerman1986-unordered}{The equality-pattern obstruction to female parity}
\begin{proof}
Counting the \(F\)-entries in an explicit unary table and reducing
the count modulo two is polynomial time.  Suppose a fixed
\(\mathcal L_\mu\) sentence defined that query.  Restrict to
\(B_n\) with \(\lvert U\rvert=n\), \(F=U\),
\(P=H=\varnothing\), and one distinguished \(c_0\).
Rename bound variables and relation variables apart in the fixed finite
syntax, and choose \(v\) larger than the total number of individual
variables appearing in it and its fixed-point argument tuples.
Take \(n,m>v+1\).

Assignments on these two structures have the same equality pattern
when exactly the same variables coincide and exactly the same variables
equal the distinguished constant.  Each atomic input formula then
agrees: equality agrees, \(F\) is always true, and \(P,H\)
are always false.  Any existential witness extends a matching pattern
on the other side: if it equals an assigned point or \(c_0\),
use that image; otherwise choose a fresh point, available because the
assignment uses at most \(v\) points and the domain has more than
\(v+1\).  This works in both directions.  Boolean operations and
quantifiers therefore preserve truth determined by such patterns.

For the syntax induction, allow any interpretations of free relation
variables that agree on matching tuple patterns, relative to the held
external parameters.  This strengthened assertion supplies the induction
hypothesis when an outer fixed-point stage is passed to an inner binder;
it does not assume that the recursive relation is an input predicate.
To include a fixed-point constructor, hold its external parameter
pattern fixed.  At stage zero its recursively defined relation is empty
on both structures.  Suppose at stage \(i\) corresponding tuples
have matching membership, determined by their equality pattern together
with those parameters and \(c_0\).  The structural argument just
given, applied to the defining matrix with these recursive atoms as
additional pattern-invariant atoms, gives matching membership at stage
\(i+1\).  Nested fixed points are handled by the same induction
on their smaller syntax.  The finitely many possible equality patterns
are the same in both domains, since both have enough fresh elements
to realize every pattern used by the fixed formula.  Thus the stage
operators on that finite set of patterns are identical.  Their
ascending iterations stabilize at the same pattern set, and least
fixed-point membership agrees.  This proves the structural claim
for arbitrary nesting, not just for first-order formulas or one
particular number of induction rounds.

A sentence has no assigned free variables.  The claim shows its truth
is the same in all \(B_n\) beyond this bound.  But in these
structures \(\lvert F\rvert=n\), and parity differs for
arbitrarily large successive sizes.  The putative defining sentence
cannot exist.  This proves the explicit example of the source's
imported Fact~2.2 without dropping its constant or adding an order
\cite[sec.~2]{Immerman1986}.
\end{proof}
\end{annexitem}
